% status: SKELETON
% Chapter 4 of the second edition. Plan: docs/STRUCTURE.md. Rules: AUTHORING.md.
\chapter{Prefill and Decode: Two Workloads, One Model}\label{ch:prefill-decode}

\begin{ChapterIntent}
Processing a prompt and generating a reply use the same weights and
almost nothing else in common. Prefill multiplies the weights by a matrix of
prompt tokens and is limited by compute; decode multiplies them by one vector
per sequence and is limited by memory bandwidth. This chapter measures both
phases on one machine, derives time to first token and time per output token
from the roofline, and shows why each phase wants a different batch size, a
different kernel, a different scheduler and --- at scale --- different
hardware. The tension between them drives chunked prefill
(\Chref{chunked-prefill}) and disaggregated serving
(\Chref{disaggregated-serving}).
\end{ChapterIntent}

\OpeningSection{The same model reads thirty times faster than it writes}

\NapkinSection{Napkin math: time to first token and time per output token}

\section{Prefill: a matrix of tokens}

\section{Decode: one token at a time}

\section{Attention's share grows with context}

\section{Latency, defined: TTFT, TPOT, ITL and end to end}

\section{Where the phases collide}

\section{Two phases, two ideal machines}

\LedgerSection

\LabSection{time the two phases}

\HappenedSection{interference in a shared batch}

\FrontierSection{}

\ExercisesSection
